\section{Summary of Findings}
This research advances the state of trustworthy machine learning by bridging the gap between adaptive privacy preservation and Byzantine resilience in decentralized environments. The primary contribution is the development of a Multi-Agent Reinforcement Learning (MARL) framework that transforms Differential Privacy (DP) from a static constraint into a dynamic defensive mechanism. Unlike traditional aggregation methods that often fail under the combined pressure of noise injection and adversarial poisoning, this framework introduces the Dynamic Noise Optimization (DNO-RL) method. This method intelligently calibrates Gaussian noise levels by analyzing the variance and reliability of local updates, effectively masking sensitive data while simultaneously isolating anomalous Byzantine contributions. Furthermore, by optimizing the privacy-accuracy trade-off in real-time, this work provides a scalable solution that maintains model integrity and convergence speed even under sophisticated label-flipping and gradient-manipulation attacks, offering a provable DP guarantee without the catastrophic utility loss common in existing literature.
%\section{Contributions to the Field}

\section{Recommendations for Future Work}
While this framework establishes a robust foundation for adaptive security, several avenues for future research remain to be explored to enhance its practical deployment. Subsequent studies should focus on extending the MARL controller to handle extreme data heterogeneity (Non-IID), as current robust aggregators still exhibit performance degradation when local data distributions vary significantly across clients. There is also a critical need to investigate the communication overhead associated with MARL training, potentially through the integration of sparse communication techniques or model compression to ensure the framework is viable for resource-constrained IoT devices. Additionally, future work could incorporate Hardware-Assisted Security, such as Trusted Execution Environments (TEEs), to provide an extra layer of protection for the MARL policy updates themselves against server-side compromise. Finally, expanding the threat model to include stealthy backdoor attacks and evaluating the system in high-stakes real-world domains, such as robotic-assisted surgery or financial fraud detection, will be essential to validating the long-term reliability and safety of adaptive privacy-preserving collaborative learning.
